\begin{prob}{006}{}{2020年11月号 学コン5番}
\begin{enumerate}
 \item $x^2 + y^2 = 3z^2$の整数解は$(x,y,z) = (0,0,0)$ に限ることを証明せよ。
 \item $xy$平面において, $x$座標と$y$座標がともに有理数であるような点のことを「有理点」と呼ぶこととする。$\theta$を$0< \theta \leq \dfrac{\pi}{2}$を満たす定数とし, 次の条件を満たす$xy$平面上の三角形$\mr{ABC}$を考える:
 \begin{center}
 (\tf{条件})\qquad $\mr{A}$は有理点であり, $\mr{AB}$の長さは有理数である。さらに, $\angle{\mr{ACB}} = \theta$である。
 \end{center}
 三角形$\mr{ABC}$の外心を$\mr{X}$とする。次の問に答えよ。
 \begin{enumerate}
  \item $\theta = \dfrac{\pi}{3}$のとき, $\mr{X}$が有理点でないことを証明せよ。
  \item $\sin{\theta}= \dfrac{1}{\sqrt{2020}}$ のとき, $\mr{X}$が有理点となることはあるか. あるならば点$\mr{A},\mr{B},\mr{C}$の座標を挙げよ.
 \end{enumerate}
\end{enumerate}
\end{prob}
